\section{Public-only construction and the development record}
\label{sec:development}

The final candidate is the frozen mechanical baseline. It received no
question-level supervision, and the metered \texttt{dev-B} partition was never
consumed. This section reports how that candidate was built from public inputs
only, what the development partition showed about it, and what the append-only
control plane recorded about the interventions that were proposed against it.

\subsection{What the public knowledge base could be compiled into}

Before any hidden development label was released, public schema, public column
meanings, and the public HKB were transformed deterministically into semantic
artifacts. Every knowledge node received exactly one disposition:
\texttt{compile}, a defensible same-table, same-grain scalar definition;
\texttt{context\_only}, useful semantic text with one exact field target but no
safe executable object; \texttt{defer\_cross\_grain}, a definition requiring
unresolved identity, relationship, cardinality, aggregation, temporal, or
ordering semantics; or \texttt{unsupported}, missing, ambiguous, or currently
non-representable inputs. Guessing a join or an aggregation would have raised
apparent model coverage while destroying the governance claim the benchmark is
meant to test, so the rule is deliberately conservative.

\begin{table}[htbp]
\centering
\caption{HKB translation coverage across all 18 databases.}
\label{tab:hkb}
\begin{tabular}{lrr}
\toprule
Disposition & Definitions & Share \\
\midrule
Compiled & 193 & 17.7\% \\
Context only & 193 & 17.7\% \\
Deferred cross-grain & 511 & 46.9\% \\
Unsupported & 193 & 17.7\% \\
\midrule
Total & 1,090 & 100.0\% \\
\bottomrule
\end{tabular}
\end{table}

Table~\ref{tab:hkb} is the single most consequential measurement in the study,
because it determines what the governed condition could possibly have done. Fewer
than one definition in five became an executable object, and almost half were
deferred for crossing a grain the public inputs never resolve. Within the
17-database fan-out beyond the canary, the three most frequent recorded loss
codes were \texttt{cardinality\_unknown} (398), \texttt{aggregation\_unspecified}
(314), and \texttt{cross\_grain\_no\_identity} (308). Domains built from
row-local physical or sensor definitions, such as planets and solar panels,
compiled comparatively well; residential and reverse-logistics models retained
useful context but compiled no HKB definition safely under the same rules.

This measures transformation coverage and not question correctness. A deferred
definition may never be needed for any benchmark question, and a compiled
definition may still be retrieved or interpreted wrongly. What it does establish
is that the deployed model published no join path and no measure, which is the
precondition for the query-path finding in Section~\ref{sec:conditions}.

\subsection{Bounding the direct comparator}

The first end-to-end vertical slice ran one public development question before
anything scaled. The first C1 attempt passed the public-question,
database-parity, read-only, model-identity, and first-turn gates, and then its
opening schema inspection returned all 51 tables of the database and the next
model turn exceeded its budget before producing any SQL. Whole-database
inspection was replaced with deterministic query-directed search over the same
committed public schema: the model must supply a non-empty query, and the tool
returns at most four matching tables inside a 64 KiB payload. C1 through C3 share
that retrieval contract.

\begin{table}[htbp]
\centering
\caption{One question, three integration attempts. Each row is a single attempt
on one database, not a population estimate.}
\label{tab:slice}
\begin{tabular}{lrrrr}
\toprule
Stage & Terminal state & Tokens & Cost & Latency \\
\midrule
Whole-schema C1 & Budget error & 173,365 & \$1.74 & 26.0 s \\
First bounded diagnostic & Public-ID validation error & 1,585 & \$0.02 & 3.0 s \\
Reviewed bounded C1 canary & Answered & 33,445 & \$0.21 & 40.9 s \\
\bottomrule
\end{tabular}
\end{table}

The intermediate failure in Table~\ref{tab:slice} is preserved on purpose. The
first bounded attempt still returned no answer, because a generic secret
heuristic rejected canonical public foreign-key identifiers; a typed, fail-closed
validation rule fixed that boundary while keeping the credential checks. Keeping
the failed row prevents the whole change from being attributed to retrieval
alone. The four-match window itself remains unvalidated as a choice: a later
attempt on a different archeology question reached five four-table searches, cost
\$7.49, and ended in a model budget error, and reducing the window to two tables
cut that attempt to four searches and \$4.32 but still returned no answer. The
cheaper failure was no evidence for the smaller window, so the change was
reverted, and a 20-question sensitivity arm over the match cap was frozen and
then deferred before execution.

The product lesson generalizes past this comparator. Tool payload bounds are part
of agent quality: a semantically valid tool call can make a system unusable by
consuming the remaining inference budget, and a safety filter that treats every
identifier as untrusted free text can reject legitimate public metadata.

\subsection{The development baseline}

The 630-attempt direct baseline was frozen by content hash before the train-only
gold release. Its exact development intersection contains 420 attempts over 140
of the 154 development questions; 14 questions have no baseline output and were
left missing rather than backfilled. Gold conformance left 122 questions per
condition scoreable under the official scorer, and 121 under the sensitivity
scorer, which retained one fewer result that exceeded its fixed normalization
limit.

\begin{table}[htbp]
\centering
\caption{Direct-SQL development baseline, exploratory. Percentages are over each
condition's scoreable denominator.}
\label{tab:devbaseline}
\begin{tabular}{lrrrrr}
\toprule
Condition & Correct & Wrong & Refused/error & Official & Sensitivity \\
\midrule
C1 raw schema & 9 & 80 & 33 & 7.4\% & 7.4\% \\
C2 searchable raw HKB & 29 & 91 & 2 & 23.8\% & 23.1\% \\
C3 searchable exported model & 16 & 74 & 32 & 13.1\% & 11.6\% \\
\midrule
All direct conditions & 54 & 245 & 67 & 14.8\% & 14.0\% \\
\bottomrule
\end{tabular}
\end{table}

Two things in Table~\ref{tab:devbaseline} survived into the held-out result. The
ordering did: searchable raw business knowledge produced the strongest direct
baseline, and converting it into the bounded structured model gave back more than
half of that gain. And wrong answers dominated: 245 of 366 scoreable attempts
returned a wrong result rather than declining to answer. C2 also had the highest
capture-level completion rate on the 16-database frame at 97.9\%, against 75.9\%
for C1 and 74.8\% for C3, but 91 of its 122 scoreable attempts were wrong, so
completion cannot stand in for correctness.

The governed arm was budgeted and scheduled separately, over the full fixed
development schedule rather than the direct arm's represented frame. Of its 154
scheduled questions, 136 were answerable and all 136 executed: 9 correct, 93
wrong, and 34 refused or ended in a system-contract error, for 9/136 (6.6\%)
official accuracy; the sensitivity scorer retained 135 scoreable questions with
the same 9 correct and 93 wrong, for 9/135 (6.7\%). An aggregate-only alignment
to the 122 questions scoreable in all four conditions gives the governed arm 5
correct, 83 wrong, and 34 refused or system-error, or 5/122 (4.1\%), against
7.4\%, 23.8\%, and 13.1\% for C1 through C3 on that same frame. The paired
C1/C4 table on those 122 questions holds 3 correct in both, 2 correct only in C4,
6 correct only in C1, and 111 correct in neither. The governed arm's low result
is therefore not an artifact of its wider denominator. No question identity, SQL,
row value, annotation, or per-question label left custody during that alignment.

A structural diagnostic over the SQL text, carrying no question identity, no SQL
body, no result value, and no hidden annotation, found that wrong answers used
more relations on average in every direct condition: 3.16 against 2.00 in C1,
3.28 against 2.48 in C2, and 3.04 against 1.81 in C3. Thirty of 31 window-query
attempts and 25 of 28 \texttt{DISTINCT}-query attempts were wrong, while the mere
presence of a join or an aggregate separated correct from wrong hardly at all.
The governed arm ran in the same direction: parseable correct queries averaged
1.67 relations against 2.62 for wrong answers and 2.88 for system errors. These
are associations. They point at relationship, grain, and dependency handling as
the next mechanism family to test, and they establish nothing causal.

\subsection{The control plane, and what it recorded}

Optimization on the development partition is the part of an evaluation most
likely to leak, to overfit, or to be narrated favorably afterward, so it was
constrained by a control plane whose records are append-only and whose failed
experiments therefore remain in the record. Routine optimization received a
derived public view containing only the 154 development identifiers and never a
checkpoint or held-out outcome. Generation artifacts may carry generated SQL or a
governed query, latency, token and cost observations, non-secret tool failures,
and the semantic-object identifiers the system used; they may not carry gold SQL,
test cases, hidden \texttt{external\_knowledge} identifiers, expected results, or
scores, and the validator rejects protected keys recursively, so nesting a hidden
field does not bypass the boundary. Correctness arrives only in a separate
immutable score artifact whose path and expected digest are mandatory for every
scored decision, and the control plane stores component hashes plus a combined
candidate hash so a guardian receipt cannot be rebound to a different generation
or score artifact. Each iteration appends a proposal \emph{before} the system
changes, stating the affected failure class, the mechanism, the predicted
direction, the regression risk, the surface, and why the intervention is
reusable; it then appends exactly one of \texttt{KEEP}, \texttt{REVERT},
\texttt{INCONCLUSIVE}, \texttt{INVESTIGATE}, or \texttt{ARCHIVE}. Every kept
intervention declares its scope, question-specific changes are prohibited
outright, and benchmark-specific candidates may be archived as evidence but never
promoted into the final system. The objective is ordered rather than scalar, so
an accuracy increase alone cannot authorize acceptance.

Four intervention families were preregistered against the mechanism ladder:
same-grain dependency composition (E01), FK-backed grain relationships (E02),
bounded semantic descriptions (E03), and a broad HKB-context negative control
(E04). A fifth, E05, was registered later against the governed arm's
\texttt{UNKNOWN}-type result-contract failures. None was promoted.

E01 resolved inconclusive on inspection: the frozen baseline already carried 48
dependency-bearing elements, 70 executable dependency edges, and depth three, so
there was no contrast left to run.

E02 is the substantive one. It declares FK-backed relationships, which is exactly
the ingredient whose absence left the rewrite path as the only route to
cross-table access, so it is a direct test of the mechanism rather than a
candidate picked from a structural correlation. Of 1,228 public foreign keys,
1,049 pass the conservative cardinality contract, and the bounded artifact emits
91 relationships across 16 databases and 67 source topics with zero
metric-disposition changes. Its deployment verified all 16 targets with exact
readback, and its single 136-attempt development generation completed and was
frozen. It then failed its own coverage gate: 19 of the 136 attempts captured no
scoreable outcome. Fourteen were saved semantic queries rejected by the
result-type capture contract, and five were genuine transport failures that
retained no semantic query at all, so the preregistered complete-136 score could
not be produced without a new model attempt, which the no-rerun rule forbids.
E02 is recorded INCONCLUSIVE.

Because the answer to a stopped experiment is not a rerun, a no-rerun diagnostic
was applied offline to the 117 captured answers under both frozen scorers, with
no model answer regenerated. Official accuracy on that subset was 11/117 (9.4\%)
for E02 against 9/117 (7.7\%) for the frozen governed baseline on the identical
coordinates, a paired difference of $+1.7$ points, with four transitions to
correct and two regressions from correct; sensitivity gave 10/116 (8.6\%) against
9/116 (7.8\%), or $+0.9$ points. The full-frame official bounds run from 11/136
(8.1\%) if every missing outcome is wrong to 30/136 (22.1\%) if every one is
correct, and treating the 14 contract failures as failures and only the five
transport losses as unresolved still leaves an upper bound of 16/136 (11.8\%).
Bounds that wide, over that few transitions, support further study and not a
claim that E02 improved accuracy.

E05 was stopped by its own precondition before it consumed a live attempt. It
proposed declaring explicit output types on compiled semantic fields to address
the 31 \texttt{UNKNOWN}-type contract failures, and required at least 16 of those
31 attempts to have selected a compiled derived field. Measured offline on the
immutable generation records, the ceiling is 6 of 31, and 24 of the 31 select no
compiled bundle field of any kind, so no declaration on a compiled field could
reach them. E03 and E04 were never run and fell out of scope.

\subsection{One deviation, stated in full}

The untuned held-out arm was scored before E02's development execution finished,
which reverses the ordering that the later lean optimization extension specified.
E02 had already been selected and preregistered as a relationship-path contrast,
and its general compiler change predates the held-out results, so the frozen
untuned comparison is unaffected: every held-out generation completed before
scoring, and both scorers were published in one atomic run. The cost falls
entirely on optimization. Because held-out aggregates are now visible, they may
not drive a new intervention edit, a checkpoint consumption, a promotion
decision, or an optimized held-out arm, and E02 cannot be rerun or promoted into
a sealed successor. The study therefore reports a valid frozen four-condition
held-out comparison and a separate pre-specified development mechanism
experiment. It makes no claim of held-out improvement from tuning. One later
correction generalized the publication of already-normalized relationship
endpoints in response to public validator evidence; it used no held-out outcome
and changed no experimental choice.
